\documentclass[10pt,a4paper]{amsart}
\usepackage[margin=19mm]{geometry}
\usepackage{amsmath,amssymb,mathtools}
\usepackage[scr=rsfs,cal=euler]{mathalfa}
\usepackage{xfrac}
\usepackage[unicode,hidelinks,bookmarks=false]{hyperref}
\newcommand{\ie}{i.e.\ }
\newcommand{\cf}{cf.\ }
\newcommand{\resp}{respectively\ }
\newcommand{\Subsection}{\subsection}
\providecommand{\nobreakdash}{\nobreak\hbox{-}}
\setlength{\emergencystretch}{2em}
\multlinegap0pt
\usepackage{bm}
\usepackage{bbm}
\usepackage{dsfont}
\usepackage{xspace}
\usepackage{cancel}
\usepackage{mathtools}
\usepackage{amsthm}
\makeatletter
\@ifundefined{theorem}{\newtheorem{theorem}{Theorem}}{}
\@ifundefined{lemma}{\newtheorem{lemma}[theorem]{Lemma}}{}
\makeatother
\title[Parametric Scaling Law of Tuning Bias in Conformal Predictio]{Parametric Scaling Law of Tuning Bias in Conformal Prediction}
\author{Hao Zeng, Kangdao Liu, Bingyi Jing, Hongxin Wei}
\date{Source: arXiv:2502.03023v2 (CC BY 4.0)}
\begin{document}
\maketitle

\noindent\emph{Source and licence.} Parametric Scaling Law of Tuning Bias in Conformal Prediction. Version arXiv:2502.03023v2 (CC BY 4.0), distributed under the Creative Commons Attribution 4.0 International licence, as recorded in the arXiv licence metadata for this exact version and shown on the abstract page https://arxiv.org/abs/2502.03023v2. Full rights notice: licenses/arxiv-2502.03023-CC-BY-4.0.txt.\par\smallskip

\noindent\textit{Editorial scope.} Counted unit: the Lemma whose source label is \texttt{lemma:order\_preserving\_vector\_scaling} in the article (source0.tex lines 1489--1492, source label \texttt{lemma:order\_preserving\_vector\_scaling}), delivered together with the article's own complete original proof of it (source0.tex lines 1494--1521, one \texttt{proof} environment of 28 lines). No mathematics is written, repaired or re-derived: statement and proof are the article's own text line for line. Required context retained uncounted: none. Every definition, display and label of those lines is retained unchanged. Omitted from this package and not redistributed: 1--1488, 1493--1493, 1522--1564. Nothing inside a retained excerpt is omitted or altered: every retained body is delivered exactly as the article's lines are written, so no omission receipt of any kind accompanies this package. Citations: the retained excerpts carry no citation command at all, so no bibliography entry of the article is transcribed and no reference is redistributed. Reference adaptation: none was needed and none was made; every reference inside the delivered unit resolves inside it, so no unresolved reference or citation is produced. Printed slips kept literal: no printed word, symbol, hypothesis or number of the counted statement or of its proof was altered. Layout and class: the article's own class is \texttt{article} with packages including cite, microtype, graphicx, booktabs, hyperref, float, subcaption, amsmath, amssymb, mathtools, amsthm, dsfont, icml2025, cleveref; the wrapper is the lane's pinned \texttt{amsart} editorial wrapper at 10pt on A4, which restates the theorem-like environments the retained text needs and adds the article's own macro definitions that the retained text needs (none), with extra wrapper packages (bm, bbm, dsfont, xspace, cancel, mathtools). The delivered numbering is the unit's own. Assets: no retained span contains a figure environment, a graphics-inclusion command, a PDF-page inclusion, a table or a TikZ picture; no image file is copied into this package, the package declares no asset dependency and no image file is redistributed with it. Licence: version arXiv:2502.03023v2 of this article is distributed under the Creative Commons Attribution 4.0 International licence, as recorded in the arXiv licence metadata for this exact version and shown on the abstract page https://arxiv.org/abs/2502.03023v2. A full rights notice is delivered at licenses/arxiv-2502.03023-CC-BY-4.0.txt. Characters: every retained excerpt is pure ASCII, so the delivered engine, XeLaTeX with its default Latin Modern text font, reports no missing character.
\begin{lemma}
    \label{lemma:order_preserving_vector_scaling}
    Vector scaling is order-preserving if and only if \(W_j = W_{j'} >0\) for all \(j, j'\in [K]\) and \(b_j = b_{j'}\) for all \(j, j'\in [K]\).
\end{lemma}

\begin{proof}
    \textbf{\((\Leftarrow)\)}
    It is a direct result of the definition of vector scaling.

    \textbf{\((\Rightarrow)\)}
    If vector scaling is order-preserving, then
    \begin{equation*}
        \forall j, j'\in [K], f_{j}(\boldsymbol{x}) \leq f_{j'}(\boldsymbol{x}) \iff W_j f_{j}(\boldsymbol{x}) + b_j \leq W_{j'} f_{j'}(\boldsymbol{x}) + b_{j'}.
    \end{equation*}
    We first exclude the case when \(W_j <0\) for any \(j\in [K]\):
    For any \(W_j <0\), and fixed \(b_j, b_{j'}\), we could find small enough \(f_{j}(\boldsymbol{x})\), such that \(f_{j}(\boldsymbol{x}) < f_{j'}(\boldsymbol{x})\) and \(W_j f_{j}(\boldsymbol{x}) + b_j > W_{j'} f_{j'}(\boldsymbol{x}) + b_{j'}\), which contradicts the order-preserving property.
    Therefore, we only consider the case when \(W_j, W_{j'} >0\) for all \(j, j'\in [K]\). 
    Without loss of generality, we assume \(f_{j}(\boldsymbol{x}) \leq f_{j'}(\boldsymbol{x})\) and \(W_j < W_{j'}\), and let \(f_{j'}(\boldsymbol{x}) = f_{j}(\boldsymbol{x}) + \epsilon\) with \(\epsilon \geq 0\), then we have
    \begin{equation*}
        W_j f_{j}(\boldsymbol{x}) + b_j \leq W_{j'} f_{j'}(\boldsymbol{x}) + b_{j'} \iff 
        (W_j - W_{j'})f_{j}(\boldsymbol{x})  \leq b_{j'} - b_j+ W_{j'} \epsilon
    \end{equation*}
    Let \(\varepsilon \to 0\), we have
    \begin{equation*}
        (W_j - W_{j'})f_{j}(\boldsymbol{x})  \leq b_{j'} - b_j
    \end{equation*}
    hold for any \(f_{j}(\boldsymbol{x})\). We could find a small enough \(f_{j}(\boldsymbol{x})\), such that
    \begin{equation*}
        (W_j - W_{j'})f_{j}(\boldsymbol{x})  > b_{j'} - b_j
    \end{equation*}
    This contradicts the order-preserving property.
    Therefore, we must have \(W_j = W_{j'} >0\) for all \(j, j'\in [K]\) and \(b_j = b_{j'}\) for all \(j, j'\in [K]\).
\end{proof}

\end{document}
